
% =============================================================================
% Stage24 Manuscript Integration
% Insert into an IEEE manuscript as appropriate.
% =============================================================================

\subsection{Bidirectional Cross-Dataset Generalization}

Stage24 evaluated transfer in two separately reported directions. The
IDS2018-to-CICIDS2017 direction used 2,830,743 effective
CICIDS2017 flows with attack prevalence 0.196996, whereas the
reciprocal CICIDS2017-to-IDS2018 direction used the frozen K79-clean IDS2018
Feb-28 population of 593,780 flows with attack prevalence
0.104847.

For IDS2018-to-CICIDS2017, bridge62 achieved PR-AUC
0.667483 and ROC-AUC
0.733946. Bridge70 under the
published semantics achieved PR-AUC
0.663981 and ROC-AUC
0.741891, while the flag-corrected
bridge70 representation achieved PR-AUC
0.656252 and ROC-AUC
0.744083.

The reciprocal transfer was substantially weaker. CICIDS2017-trained bridge62
reached PR-AUC 0.108176 and ROC-AUC
0.525167; bridge70 reached PR-AUC
0.107419 and ROC-AUC
0.525302. The IDS2018 Feb-28 PR-AUC
prevalence anchor was 0.104847. The two transfer directions
are therefore retained separately rather than averaged.

\subsection{Extractor-Semantic Sensitivity}

Within the primary bridge70 evaluation, replacing the published CICIDS2017
aggregate-flag interpretation with the frozen corrected mapping changed PR-AUC
by -0.007729
(95\% CI [-0.007921, -0.007534]),
ROC-AUC by +0.002192
([+0.002139, +0.002246]), and Brier
score by +0.005087
([+0.005023, +0.005152]).
All three intervals exclude zero, demonstrating measurable sensitivity to
extractor/serialization semantics.

\subsection{Operating-Point Transfer}

CICIDS2017 Thursday source-validation thresholds transferred poorly to IDS2018
Feb-28. For bridge62, standard, balanced, and security recall were
0.000386,
0.000610, and
0.001349,
respectively. Corresponding bridge70 recall values were
0.000386,
0.000610, and
0.000819.
No target threshold search was permitted.

\subsection{Discussion}

The Stage24 results demonstrate that cross-dataset IDS portability can be
strongly directional even when feature names appear compatible across benchmark
families. Successful IDS2018-to-CICIDS2017 transfer did not imply reciprocal
CICIDS2017-to-IDS2018 generalization. The bridge ablation also shows that adding
aggregate flag-count fields is not uniformly beneficial, while the serialization
audit demonstrates that semantically misaligned extractor outputs can
materially change reported cross-dataset results. Ranking discrimination,
probability calibration, and operating-point transfer should therefore be
evaluated as separate properties.

\subsection{Limitations and Threats to Validity}

The experiment covers two related CIC benchmark families and one frozen model
strategy per direction, so the observed magnitude of asymmetry should not be
generalized to unrelated operational networks or all model families. Target
prevalence differs between directions, requiring PR-AUC to be interpreted
together with prevalence, normalized PR-AUC, and ROC-AUC. Two preregistered
GROUNDED\_S4 cells were administratively cancelled because exact durable
physical-row membership could not be reconstructed without a new post-freeze
heuristic; no substitute population was used. Finally, paired row-bootstrap
intervals characterize uncertainty for the frozen target populations and do
not represent variation across independently collected organizations or
capture campaigns.

\subsection{Stage24 Contribution}

A validation-safe bidirectional cross-dataset audit revealed pronounced
transfer asymmetry between CSE-CIC-IDS2018 and CICIDS2017, quantified the
contribution of a common 62-feature bridge versus aggregate TCP-flag fields,
demonstrated statistically resolved sensitivity to CICIDS2017 flag
serialization semantics, and showed that source-selected operating thresholds
can fail catastrophically under reciprocal transfer without target-guided
adaptation.
